\section{Conclusion}
\label{sec:conclusion}

Based on the analysis of test data and performance evaluation, the
following critical conclusions are drawn regarding the proposed
\textbf{PredCLOCK} system:

\begin{enumerate}
  \item \textbf{Proven Write Efficiency with Conservative Limits:}
    The implementation of \textit{PredCLOCK} successfully validated
    the hypothesis that a predictive model can reduce write
    operations (promotions) to main memory. A consistent reduction of
    \textbf{8--10\%} was achieved across all cache size variations
    compared to the Standard CLOCK baseline.
    While this reduction helps extend the endurance of modern storage
    devices (such as NVM or SSDs), it falls significantly short of
    the theoretical optimal bound of \textbf{43\%} achieved by the
    \textit{Offline CLOCK} oracle. This gap indicates that the model,
    bounded by local feature visibility, adopts a conservative
    decision-making process to preserve cache hits.

  \item \textbf{Criticality of Handling Class Imbalance:}
    The parameter tuning phase revealed that the extreme class
    imbalance in Zipfian workloads (where "one-hit wonders" dominate
    "reused" objects by approx. 80:20) is the primary challenge for
    ML-based caching.
    \begin{itemize}
      \item The \textbf{Aggressive Model} (Unweighted Loss) failed
        systemically, yielding a 2\% regression in Miss Ratio
        (Service Degradation) due to bias toward the majority class.
      \item The \textbf{Conservative Model} (Weighted Loss 2:1)
        successfully mitigated this risk, mimicking the safety of
        Standard CLOCK while filtering obvious waste. This confirms
        that cost-sensitive learning is mandatory for cache
        replacement policies.
    \end{itemize}

  \item \textbf{Non-Linear Scalability and Over-Selectivity:}
    There is a distinct non-linear performance pattern relative to
    cache size. In high-pressure environments (1\% -- 20\% capacity),
    \textit{PredCLOCK} provides a "double benefit" of reduced miss
    ratios and reduced promotions. However, a performance regression
    (+0.15\% Miss Ratio) occurs at a cache size of 40\%. At this
    capacity, the static decision threshold (0.8) becomes
    \textbf{over-selective}, discarding borderline objects that
    should have been stored given the ample available space. This
    confirms that static ML thresholds lack the flexibility required
    for varying memory pressures.

  \item \textbf{Superiority of Oracle Knowledge:}
    The dominance of the \textit{Optimal Zipf Promotion} strategy
    over \textit{PredCLOCK} across all metrics confirms that in
    Zipfian distributions, global knowledge regarding object
    popularity is far more valuable than the local access patterns
    (recency and short-term frequency) utilized by \textit{PredCLOCK}
    .

  \item \textbf{Viability of Hybrid Architecture:}
    Technically, this research demonstrated the viability of a hybrid
    simulation framework. The integration of Python-trained models
    into a high-performance C++ environment via \textbf{ONNX (Open
    Neural Network Exchange)} proved stable and efficient, bridging
    the gap between modern deep learning tools and operating system
    simulations.
\end{enumerate}
